% =============================================================================
%  Lab 6 - Register File
% =============================================================================

\chapter{Lab 6 --- Register File}

\section{Objective}
% Two or three sentences, your own words.

\section{Specification}

\begin{table}[htbp]
\centering
\caption{Compressed table (Type C) for the \texttt{RegFile} write port.}
\label{tab:lab06-truth}
\begin{tabular}{ccl||l}
\toprule
\# & \texttt{Strb} & \texttt{Dsel} & effect at the rising edge \\
\midrule
0 & 0 & --- & no write \\
1 & 1 & \texttt{"00000"} & \textbf{no write} --- register 0 is hardwired \\
2 & 1 & $k \neq 0$ & $R_k \leftarrow \texttt{Din}$ \\
\bottomrule
\end{tabular}
\end{table}

Reads are combinational and always valid:
$\texttt{Aout} = R[\texttt{Asel}]$, $\texttt{Bout} = R[\texttt{Bsel}]$,
with $R_0 \equiv 0$.

\subsection*{Equations}

\begin{align*}
\text{address minterm} \quad & m_k(S) = 1 \text{ iff } S = k \\
\text{write enable}    \quad & \mathit{We}_k = \texttt{Strb} \cdot m_k(\texttt{Dsel}) \cdot [k \neq 0] \\
\text{storage}         \quad & R_k^{+} = \mathit{We}_k \cdot \texttt{Din} + \overline{\mathit{We}_k} \cdot R_k \\
\text{read port A}     \quad & \texttt{Aout} = \sum_{k} m_k(\texttt{Asel}) \cdot R_k
\end{align*}

% ONE SENTENCE, required: circle or name the term that makes register 0
% unwritable, and say what "hardwired to zero" means in hardware terms.

\section{Design}
% Answer all three:
%   (a) A register file is three components you already built. Name them, say
%       which lab each came from, and point at the lines of your code where
%       each one lives.
%   (b) Which read-during-write behaviour did you build - READ-OLD or READ-NEW?
%       How does your testbench demonstrate it? Both are valid designs; the
%       failure would be not knowing which one you have.
%   (c) Why are the read ports OUTSIDE the clocked process? What would break in
%       Lab 12 if they were inside it?

\section{Implementation}

\vhdlfile{code/lab06/RegFile.vhd}{\texttt{RegFile} --- 32 $\times$ 32, two read ports, one write port}

\section{Verification}
% Vectors, then the coverage argument.
%
% You cannot test 32 registers x 2^32 values. The argument is ADDRESS-SPACE
% BOUNDARIES plus the special case:
%     register 0 (the special one), register 31 (the top), one in the middle
% State that in your own words.

\begin{table}[htbp]
\centering
\caption{Test vectors.}
\label{tab:lab06-vectors}
\begin{tabular}{clllp{4.6cm}}
\toprule
\# & action & address & expected & what wrong design this catches \\
\midrule
1 & & & & \\
2 & & & & \\
3 & & & & \\
\bottomrule
\end{tabular}
\end{table}

\section{Results}
% Evidence: E1, E2, E4, E6, E7.

\figplaceholder{the waveform showing an attempted write to register 0, and
register 0 still reading zero afterwards}
% Required sentence: state that Strb was asserted with Dsel = 0 and that the
% stored value did not change, and name the mechanism that prevented it.

\figplaceholder{the RTL Viewer, showing the storage array and the read multiplexers}

% E4: the register count. THIS IS THE FIRST DESIGN THAT SHOULD HAVE REGISTERS.
%     Expect around 1024 (32 x 32). Quote your number and explain it. A much
%     lower count usually means the Fitter inferred block memory instead of
%     flip-flops - say which happened to you and how you can tell.
% E6: Fmax, and one sentence on where the path runs.

\section{Discussion}
% REQUIRED. One bug, how you found it, what it cost.
%
% Then answer this explicitly:
%     Register 0 holds nothing and can never be changed, yet it occupies one of
%     only 32 register addresses. Why is that a good trade?
% Give at least one instruction that becomes unnecessary because $0 exists, and
% point at where the Lab 12 test program relies on it.
